\begin{exercise}{151}
Let $\mathcal{C}$ be a cocomplete $\infty$-category and let $F\colon\Delta^n\to\mathcal{C}$ be a functor. Calculate
the colimit of the restriction $\partial\Delta^n\to\Delta^n\to\mathcal{C}$.

\begin{proof}
Write $x_i=F(i)$ and
$L_n(F)=\operatorname*{colim}_{\partial\Delta^n}F$.
The base cases are
\[
 L_0(F)\simeq\varnothing_{\mathcal{C}},\qquad
 L_1(F)\simeq x_0\amalg x_1,
\]
where $\varnothing_{\mathcal{C}}$ denotes an initial object.
For $n\ge2$, put
$D=\Delta^{\{0,\ldots,n-1\}}$, $E=\partial D$, and $G=F|_D$.
Then the answer is the finite recursion
\[
 L_n(F)\simeq x_{n-1}\amalg_{L_{n-1}(G)}x_n.
\]
The two structure maps from $L_{n-1}(G)$ are induced by the cones
\[
 E\star\{*\}\longrightarrow D\xrightarrow{\ F\ }\mathcal{C},
 \qquad
 E\star\{n\}\hookrightarrow\Delta^n\xrightarrow{\ F\ }\mathcal{C}.
\]
The first map fixes $E$ and sends the new vertex $*$ to $n-1$.
Thus the maps in the pushout include the coherent cone data supplied
by $F$.

Indeed, there is an equality of subcomplexes
\[
 \partial\Delta^n=D\amalg_E(E\star\{n\}).
\]
The restriction to $D$ has colimit $x_{n-1}$ by
Corollary~4.3.20.
The restriction to $E\star\{n\}$ has colimit $x_n$ by
Proposition~4.3.3 and Remark~4.3.5: inclusion of the cone vertex
is right-anodyne.
Proposition~4.3.26 now gives the asserted pushout with exactly the
two restricted cones above.
Finally, $\partial\Delta^0=\varnothing$ and
$\partial\Delta^1=\{0\}\amalg\{1\}$ give the stated base cases.
\end{proof}
\end{exercise}
